%% Packages, typography and theorem environments.
%% The mathematical notation lives in notation.tex, the numerals in results.tex.

%% ------------------------------------------------------------------ packages
\usepackage[T1]{fontenc}
\usepackage[margin=1.15in]{geometry}
\usepackage{amsmath,amssymb,amsthm,mathtools}
\usepackage{booktabs,array}
\usepackage{microtype}
\usepackage{xcolor}
\definecolor{refblue}{RGB}{20,60,140}
\definecolor{todored}{RGB}{160,30,30}
\usepackage[colorlinks=true,linkcolor=refblue,citecolor=refblue,
            urlcolor=refblue,linktocpage=true,
            bookmarksnumbered=true,bookmarksdepth=2,
            pdftitle={Improved exponents for 3SUM and APSP from triangles in
              sparse lopsided graphs},
            pdfauthor={Jihoon Hyun},
            pdfkeywords={3SUM, APSP, Exact Triangle, thin matrix products,
              formalized mathematics}]{hyperref}

%% ---------------------------------------------------------------- typography
\allowdisplaybreaks
%% Lean names are set in a typewriter face and do not hyphenate.
\emergencystretch=3em

%% The parts of a statement are labeled (i), (ii), ...
\renewcommand{\theenumi}{\roman{enumi}}
\renewcommand{\labelenumi}{\textup{(\theenumi)}}

%% -------------------------------------------------------- Lean names
%% \lean{name} typesets the name of a Lean declaration in a typewriter face.
%% It is used every time a result is machine-checked.  The name is set by the
%% machinery of the url package (loaded by hyperref), so underscores, primes
%% and dots are safe and a long name may break after an underscore or a dot.
%% Do not use it in section titles or captions.  \leanraw is the unbreakable
%% form, \texttt{\detokenize{...}}.
\DeclareUrlCommand\leanurl{\urlstyle{tt}%
  \def\UrlBreaks{\do\.\do\_\do\/}\def\UrlBigBreaks{}%
  \def\UrlSpecials{\do\'{\mbox{\textquotesingle}}}}
\newcommand{\lean}[1]{\leanurl{#1}}
\newcommand{\leanraw}[1]{\texttt{\detokenize{#1}}}

%% -------------------------------------------------------- status markers
%% A numbered statement with no marker in its heading is machine-checked in
%% Lean, and its text names the declarations.  The one marker below goes into
%% the heading of the other statements, all of which are in Section 9:
%%   \nf    proved here on paper, not machine-checked --- never "unproved".
%% See the Formalization paragraph of the introduction.
\newcommand{\nf}{\textup{not formalized}}

%% A visible note for the parts of the draft that are still open.
\newcommand{\TODO}[1]{\textcolor{todored}{[\textsc{todo}: #1]}}

%% ------------------------------------------------------ theorem environments
\theoremstyle{plain}
\newtheorem{theorem}{Theorem}[section]
\numberwithin{equation}{section}
\numberwithin{table}{section}
\newtheorem{proposition}[theorem]{Proposition}
\newtheorem{lemma}[theorem]{Lemma}
\newtheorem{corollary}[theorem]{Corollary}

\theoremstyle{definition}
\newtheorem{definition}[theorem]{Definition}
\newtheorem{example}[theorem]{Example}

\theoremstyle{plain}
\newtheorem{question}[theorem]{Question}

\theoremstyle{remark}
\newtheorem{remark}[theorem]{Remark}
